\section{Data Exploration: Manipulation, Derivation, Consolidation and Preparation}\label{sec:data}

This section follows the four preparation steps named in the brief: \textbf{data modification}
(repairing and standardising values), \textbf{data derivation} (new variables from existing ones),
\textbf{data deduction} (removing redundant records and inferring the right treatment of ambiguous
ones) and \textbf{data reduction} (condensing high-cardinality text into analysable categories),
followed by \textbf{consolidation} into one analysis-ready profile per dataset and the exploratory
findings that shaped the analysis.

\subsection{Consolidated data-quality audit}

SAS describes the files as public, self-reported or masked sample data that may contain mistyped
entries and outliers. We therefore audited every file before analysis (Table~\ref{tab:quality}).
The guiding rule was \emph{repair what is recoverable, flag what is uncertain, delete only what is
provably redundant}.

\begin{table}[H]\centering\scriptsize
\caption{Data-quality audit and actions for the four SAS files.}\label{tab:quality}
\begin{tabularx}{\linewidth}{@{}l r r L{3.2cm} L{2.6cm} X@{}}
\toprule
File & Raw & Used & Duplicates & Missing & Other issues and actions \\
\midrule
Analytics Jobs & 15,841 & 14,840 & 1,001 exact duplicate rows removed & job\_description 23\%, job\_type 78\% & 12,907 truncated tags repaired; 814 job\_type variants merged; 2,347 multi-city rows split; text descriptions truncated at source \\
DataScience Jobs & 1,602 & 1,602 & 0 duplicates & 0 missing & salaries stored as text (`7.8L') parsed to lakh; 39 high-salary outliers kept (senior roles) \\
JDS Skill Traits & 139 & 139 & 2 reused IDs kept (different people) & 0 missing & all scores within 1--5; ceiling: 58\% at 5.0 (storytelling); column name fixed \\
SDS Personality Traits & 161 & 161 & 9 reused IDs kept (7 with different outcomes) & 0 missing & two malformed column names fixed (leading space, spaces in target name) \\
\bottomrule
\end{tabularx}

\end{table}

\subsection{Data modification: repairing and standardising values}

\textbf{Missing and inconsistent values.} \texttt{job\_type} is 78\% missing and, where present,
spelt five ways (Analytics, analytics, ANALYTICS, analytic, Analytic). We merged the 814 variants
but do not use the variable for segmentation. \texttt{job\_description} is missing in 23\% of rows,
and 11,469 descriptions are cut off at about 100 characters and end in ``\ldots'' (a property of the
source, not of our processing). The recruiter-provided \texttt{key\_skills} field is almost complete
(one missing row), so it became the primary skill signal.

\textbf{Skill tags.} The 74,884 comma-separated tags contained 12,907 entries truncated by the
export (``Rest\ldots'', ``Oracle Workflow\ldots''), which we repaired by stripping the ellipsis rather
than discarding them; 912 empty or ``\ldots''-only tags were dropped, and 28 case duplicates within a
posting removed (Figure~\ref{fig:fig_cleaning.png}).

\textbf{Headers and formats.} SDS headers contained a leading space (`` extraversion'') and spaces
inside the target name, and JDS used a hyphen (\texttt{maths-stats\_skills}); all were normalised.
DataScience Jobs salaries were stored as text (``7.8L'') and converted to lakh per year.

\fig{fig_cleaning.png}{Modification and deduction actions on Analytics Jobs (left) and missing values
after de-duplication (right).}

\subsection{Data derivation: new variables}

Text ranges were converted to numbers so they can be modelled
(Table~\ref{tab:derived}). Salary is only available as six ordered bands, so it is analysed
\emph{ordinally} rather than as an exact amount (Table~\ref{tab:bands}). The median posting asks
for 3--7 years of experience, and 2,347 postings list several cities, which were split into a list.

\begin{table}[H]\centering\small
\caption{Variables derived during preparation.}\label{tab:derived}
\begin{tabularx}{\linewidth}{@{}L{3.6cm}L{4.2cm}Y@{}}
\toprule
Derived variable & From & Rule \\
\midrule
\texttt{exp\_min}, \texttt{exp\_max} & \texttt{experience} (``5-10 yrs'') & numeric range bounds \\
salary band (ordinal), \texttt{sal\_mid\_lakh} & \texttt{salary} (``10to15'') & 6 ordered levels; midpoint for description only \\
\texttt{locations}, \texttt{multi\_city} & \texttt{location} & split on commas; flag if $>1$ city \\
\texttt{family} & \texttt{job\_desig} & 10 role families from ordered keyword rules (Section~\ref{sec:families}) \\
canonical skills & \texttt{key\_skills} & ESCO/O*NET label or data-derived label (Section~\ref{sec:skills}) \\
JDS skill-area flags & canonical skills & keyword patterns per JDS area (used for triangulation) \\
\texttt{avg/min/max\_lakh} & DataScience Jobs salaries (``7.8L'') & numeric lakh per year \\
\bottomrule
\end{tabularx}
\end{table}

\begin{table}[H]\centering\small
\caption{Distribution of the salary band (Analytics Jobs, 14,840 postings).}\label{tab:bands}
\begin{tabular}{@{}l r r r r r r@{}}
\toprule
Salary band (lakh/yr) & 0--3 & 3--6 & 6--10 & 10--15 & 15--25 & 25--50 \\
\midrule
Postings & 2,565 & 2,032 & 2,875 & 3,503 & 2,873 & 992 \\
Share & 17.3\% & 13.7\% & 19.4\% & 23.6\% & 19.4\% & 6.7\% \\
\bottomrule
\end{tabular}

\end{table}

\subsection{Data deduction: redundant records and ambiguous cases}

\textbf{Exact duplicates.} 1,001 Analytics Jobs rows were copies of another posting in every column
except the row number. They were removed (15,841 $\rightarrow$ 14,840); keeping them would
double-count demand. DataScience Jobs, JDS and SDS contain no fully identical rows.

\textbf{Re-used IDs, inferred to be different people.} JDS has 2 and SDS 9 IDs that appear twice.
The paired rows have entirely different scores, and in SDS 7 of the 9 pairs have different outcomes.
We deduce that these are masked or re-used IDs, not repeated people, and keep all rows; deleting
them would discard 2.9\% (JDS) and 11.2\% (SDS) of the employee data.

\textbf{Outliers, kept with reasons.} 39 DataScience Jobs rows are high-salary outliers by the IQR
rule (e.g.\ Senior Data Scientist at Flipkart, 82 lakh). They are plausible senior roles with few
openings, so they are kept and salaries are summarised with medians and job-weighted averages.
Low AI \& ML scores (18 in JDS) and 5 agreeableness values (SDS) are valid scale values; rank-based
methods are robust to them.

\textbf{A method deduced to be unreliable.} Meaning-based extraction of skills from the truncated
descriptions was tested and rejected (Section~\ref{sec:skills}).

\subsection{Data reduction: from high-cardinality text to analysable categories}\label{sec:reduction}

Two fields are far too fine-grained to analyse directly: 10,096 distinct job titles and 8,460 distinct
skill tags. Table~\ref{tab:reduction} summarises how each was reduced while keeping the information
needed for the research questions.

\begin{table}[H]\centering\small
\caption{Data reduction steps.}\label{tab:reduction}
\begin{tabularx}{\linewidth}{@{}L{3.2cm}L{2.7cm}L{2.7cm}Y@{}}
\toprule
Field & Before & After & Method\\
\midrule
Job titles & 10,096 distinct & 10 role families & ordered keyword rules (Section~\ref{sec:families})\\
Skill tags & 8,460 distinct (73,936 uses) & 6,604 canonical skills & taxonomy mapping + merging near-duplicate spellings (Section~\ref{sec:skills})\\
Canonical skills & 6,604 & top 30--40 per model & demand ranking; rarer skills kept for description only\\
Canonical skills & 6,604 & 5 JDS skill areas & keyword patterns aligned with the JDS definitions (triangulation, RQ4)\\
Salary & 6 text bands & 1 ordinal variable & ordered levels; midpoint used only for description\\
DataScience Jobs & 1,602 company $\times$ title rows & 10 titles & job-weighted aggregation\\
\bottomrule
\end{tabularx}
\end{table}

\subsubsection*{Skill normalisation: from 8,460 free-text tags to comparable skills}\label{sec:skills}

Recruiters write the same skill many ways (``Machine Learning'', ``machine learning'', ``ML''). To
count and compare skills, the 8,460 distinct tags were mapped onto the international ESCO v1.2.1
and O*NET 30.3 taxonomies (22,758 skills) in a fixed order: exact name or acronym match;
a known skill name contained in the tag (``Core Java'' $\rightarrow$ Java); and finally the nearest
skill by sentence-embedding similarity, accepted only above 0.85 (0.90 for one-word tags). A manual
spot-check found and fixed systematic errors, for example ``HTML'' $\rightarrow$ PHP, ``Core Java''
$\rightarrow$ ``core apples'' and ``Analytical'' $\rightarrow$ analytical chemistry; four such
mappings are corrected by an explicit, documented override list.

\fig{fig_tag_mapping.png}{Mapping of recruiter tags to ESCO/O*NET (left) and similarity of
non-exact tags to their nearest skill (right).}

Tags that could not be mapped confidently (42.3\% of tag uses) are not discarded. They become
\emph{data-derived skills}, and near-duplicate spellings are merged when two tags are
$\geq 0.92$ similar (1.3\% of uses). As a result, \textbf{every tag use resolves to a canonical skill}
(Figure~\ref{fig:fig_canonical.png}), giving 6,604 distinct skills.

\fig[0.82\linewidth]{fig_canonical.png}{Every recruiter tag resolves to a canonical skill.}

\textbf{A method we tested and rejected.} We also tried extracting skills by meaning from the job
descriptions. Calibrated against the recruiters' own tags, it reached at best F1 = 6.8\%
(precision 6.3\%, recall 7.3\%; Figure~\ref{fig:fig_text_calibration.png}), and it produced
matches such as ``Hunting role'' $\rightarrow$ animal hunting. Because the descriptions are
truncated at the source, only exact skill names found in the text are used.

\fig[0.72\linewidth]{fig_text_calibration.png}{Calibrating meaning-based text matching against
recruiter tags: agreement never exceeds F1 $\approx$ 7\%, so it was excluded.}

\subsubsection*{Role families}\label{sec:families}

Job titles are highly varied (10,096 distinct titles, often padded with ``Urgent'', ``Walk-in'' or city
names). An ordered set of readable keyword rules assigns each title to one of ten families; generic
``Analyst'' falls back to Analytics/BI. Four families form the \textbf{data roles} (5,138 postings):
Data Science/ML, Data Engineering, Analytics/BI and Business Analysis/Product. Only 9.4\% of postings
remain ``Other'' (Figure~\ref{fig:fig_role_families.png}).

\fig{fig_role_families.png}{Postings and median salary band by role family; data roles in blue.}

\subsection{Consolidation: dataset profiles and exploratory findings}

\textbf{DataScience Jobs.} The 1,602 rows aggregate 93,005 openings across 642 companies and 10
standard titles; the row median of the average salary is 11.9 lakh and the median minimum experience
is 2 years.

\textbf{JDS and SDS.}

Both files are complete (no missing values), and both outcomes are balanced
(JDS: 73 high / 66 low hike; SDS: 85 high / 76 low success), so the majority-class baseline
accuracy is 52.5\% and 52.8\%. Table~\ref{tab:desc} summarises the variables.

\begin{table}[H]\centering\small
\caption{Descriptive statistics with group means by outcome.}\label{tab:desc}
\begin{subtable}{\linewidth}\centering
\caption{JDS Skill Traits (1--5 scale; outcome: salary hike).}\begin{tabular}{@{}l r r r r r r r@{}}
\toprule
Variable & Mean & SD & Min & Median & Max & Mean (high) & Mean (low) \\
\midrule
Big data & 3.85 & 0.85 & 2.3 & 3.8 & 5 & 3.94 & 3.75 \\
Maths \& stats & 4.29 & 0.84 & 2.2 & 4.6 & 5 & 4.71 & 3.83 \\
Coding & 4.27 & 0.89 & 2.2 & 4.6 & 5 & 4.64 & 3.85 \\
AI \& ML & 4.57 & 0.67 & 2.2 & 4.9 & 5 & 4.82 & 4.28 \\
Dashboards \& storytelling & 4.36 & 0.93 & 2.3 & 5 & 5 & 4.85 & 3.81 \\
\bottomrule
\end{tabular}

\end{subtable}\\[6pt]
\begin{subtable}{\linewidth}\centering
\caption{SDS Personality Traits (normalised scores; outcome: success).}\begin{tabular}{@{}l r r r r r r r@{}}
\toprule
Variable & Mean & SD & Min & Median & Max & Mean (high) & Mean (low) \\
\midrule
Neuroticism & 36.19 & 11.27 & 17 & 34 & 68 & 36.13 & 36.26 \\
Extraversion & 43.20 & 12.13 & 17 & 45 & 67 & 48.86 & 36.88 \\
Openness & 41.33 & 11.32 & 18 & 44 & 65 & 48.49 & 33.32 \\
Agreeableness & 44.60 & 11.29 & 17 & 46 & 68 & 47.72 & 41.12 \\
Conscientiousness & 45.21 & 13.21 & 18 & 49 & 66 & 53.68 & 35.74 \\
\bottomrule
\end{tabular}

\end{subtable}
\end{table}

\textbf{Exploratory findings.}
\begin{itemize}
  \item \textbf{Ceiling effect (JDS).} 58\% of juniors score the maximum 5.0 in dashboards \&
        storytelling, 49\% in AI \& ML, 45\% in coding and 41\% in maths-stats
        (Figure~\ref{fig:fig_jds_dist.png}). The scale cannot separate the strongest juniors, so
        effects are probably \emph{under}-estimated.
  \item \textbf{Group differences are visible before modelling}: high-hike juniors score higher on
        four of five skills and high-success seniors on four of five traits (Table~\ref{tab:desc}),
        motivating the formal tests in Section~\ref{sec:rq23}.
  \item \textbf{Multicollinearity} is low (maximum variance inflation factor 1.45 in JDS, 1.44 in SDS).
\end{itemize}

\fig{fig_jds_dist.png}{JDS: skill scores by salary-hike outcome. High-hike juniors cluster at the
5.0 ceiling.}
